\begin{lemma}
\label{lemma-weak-bourbaki-pre}
Let $R \to S$ be a ring map of finite presentation.
Let $N$ be a finitely presented $S$-module
which is flat as an $R$-module. Let $M$ be an $R$-module.
Let $\mathfrak q$ be a prime of $S$ lying over $\mathfrak p \subset R$.
Then
$$
\mathfrak q \in \text{WeakAss}_S(M \otimes_R N)
\Leftrightarrow
\Big(
\mathfrak p \in \text{WeakAss}_R(M)
\text{ and }
\overline{\mathfrak q} \in \text{Ass}_{\overline{S}}(\overline{N})
\Big)
$$
Here $\overline{S} = S \otimes_R \kappa(\mathfrak p)$,
$\overline{\mathfrak q} = \mathfrak q \overline{S}$, and
$\overline{N} = N \otimes_R \kappa(\mathfrak p)$.
\end{lemma}

\begin{proof}
Pick $g \in S$ as in Lemma \ref{lemma-weak-bourbaki-pre-pre}.
Apply Proposition \ref{proposition-finite-presentation-flat-at-point}
to the morphism of schemes $\Spec(S_g) \to \Spec(R)$, the quasi-coherent
module associated to $N_g$, and the points
corresponding to the primes $\mathfrak qS_g$ and $\mathfrak p$. Translating
into algebra we obtain a commutative diagram of rings
$$
\xymatrix{
S \ar[r] & S_g \ar[r] & S' \\
& R \ar[lu] \ar[u] \ar[r] & R' \ar[u]
}
\quad\quad
\xymatrix{
\mathfrak q \ar@{-}[r] \ar@{-}[rd] &
\mathfrak qS_g \ar@{-}[d] \ar@{-}[r] & \mathfrak q' \ar@{-}[d] \\
& \mathfrak p \ar@{-}[r] & \mathfrak p'
}
$$
endowed with primes as shown, the horizontal arrows are \'etale,
and $N \otimes_S S'$ is projective as an $R'$-module. Set
$N' = N \otimes_S S'$, $M' = M \otimes_R R'$,
$\overline{S}' = S' \otimes_{R'} \kappa(\mathfrak q')$,
$\overline{\mathfrak q}' = \mathfrak q' \overline{S}'$,
and
$$
\overline{N}' = N' \otimes_{R'} \kappa(\mathfrak p') =
\overline{N} \otimes_{\overline{S}} \overline{S}'
$$
By Lemma \ref{lemma-etale-weak-assassin-up-down} we have
\begin{align*}
\text{WeakAss}_{S'}(M' \otimes_{R'} N') & =
(\Spec(S') \to \Spec(S))^{-1}\text{WeakAss}_S(M \otimes_R N) \\
\text{WeakAss}_{R'}(M') & =
(\Spec(R') \to \Spec(R))^{-1}\text{WeakAss}_R(M) \\
\text{Ass}_{\overline{S}'}(\overline{N}') & =
(\Spec(\overline{S}') \to \Spec(\overline{S}))^{-1}
\text{Ass}_{\overline{S}}(\overline{N})
\end{align*}
Use Algebra, Lemma \ref{algebra-lemma-ass-weakly-ass}
for $\overline{N}$ and $\overline{N}'$. In particular we have
\begin{align*}
\mathfrak q \in \text{WeakAss}_S(M \otimes_R N)
& \Leftrightarrow
\mathfrak q' \in \text{WeakAss}_{S'}(M' \otimes_{R'} N') \\
\mathfrak p \in \text{WeakAss}_R(M)
& \Leftrightarrow
\mathfrak p' \in \text{WeakAss}_{R'}(M') \\
\overline{\mathfrak q} \in \text{Ass}_{\overline{S}}(\overline{N})
& \Leftrightarrow
\overline{\mathfrak q}' \in \text{WeakAss}_{\overline{S}'}(\overline{N}')
\end{align*}
Our careful choice of $g$ and the formula for
$\text{Ass}_{\overline{S}'}(\overline{N}')$ above shows that
\begin{equation}
\label{equation-key-observation}
\text{if }\mathfrak r' \in \text{Ass}_{\overline{S}'}(\overline{N}')
\text{ lies over }\mathfrak r \subset \overline{S}\text{ then }
\mathfrak r \subset \overline{\mathfrak q}
\end{equation}
This will be a key observation later in the proof. We will use
the characterization of weakly associated primes given in
Algebra, Lemma \ref{algebra-lemma-weakly-ass-local} without further mention.

\medskip\noindent
Suppose that
$\overline{\mathfrak q} \not \in \text{Ass}_{\overline{S}}(\overline{N})$.
Then
$\overline{\mathfrak q}' \not \in \text{Ass}_{\overline{S}'}(\overline{N}')$.
By
Algebra, Lemmas \ref{algebra-lemma-ass-zero-divisors},
\ref{algebra-lemma-finite-ass}, and
\ref{algebra-lemma-silly}
there exists an element $\overline{a}' \in \overline{\mathfrak q}'$
which is not a zerodivisor on $\overline{N}'$.
After replacing $\overline{a}'$ by $\lambda \overline{a}'$ for some nonzero
$\lambda \in \kappa(\mathfrak p)$ we can find
$a' \in \mathfrak q'$ mapping to $\overline{a}'$. By
Lemma \ref{lemma-invert-universally-injective}
the map $a' : N'_{\mathfrak p'} \to N'_{\mathfrak p'}$ is
$R'_{\mathfrak p'}$-universally injective. In particular
we see that $a' : M' \otimes_{R'} N' \to M' \otimes_{R'} N'$ is
injective after localizing at $\mathfrak p'$ and hence after
localizing at $\mathfrak q'$. Clearly this implies that
$\mathfrak q' \not \in \text{WeakAss}_{S'}(M' \otimes_{R'} N')$.
We conclude that $\mathfrak q \in \text{WeakAss}_S(M \otimes_R N)$ implies
$\overline{\mathfrak q} \in \text{Ass}_{\overline{S}}(\overline{N})$.

\medskip\noindent
Assume $\mathfrak q \in \text{WeakAss}_S(M \otimes_R N)$. We want
to show $\mathfrak p \in \text{WeakAss}_S(M)$.
Let $z \in M \otimes_R N$ be an element such that $\mathfrak q$
is minimal over $J = \text{Ann}_S(z)$.
Let $f_i \in \mathfrak p$, $i \in I$ be a set of generators of the
ideal $\mathfrak p$. Since $\mathfrak q$ lies over $\mathfrak p$, for every $i$
we can choose an $n_i \geq 1$ and $g_i \in S$, $g_i \not \in \mathfrak q$
with $g_i f_i^{n_i} \in J$, i.e., $g_i f_i^{n_i} z = 0$.
Let $z' \in (M' \otimes_{R'} N')_{\mathfrak p'}$ be the image of $z$.
Observe that $z'$ is nonzero because $z$ has nonzero image in
$(M \otimes_R N)_\mathfrak q$ and because $S_\mathfrak q \to S'_{\mathfrak q'}$
is faithfully flat. We claim that $f_i^{n_i} z' = 0$.

\medskip\noindent
Proof of the claim: Let $g'_i \in S'$ be the image of $g_i$.
By the key observation (\ref{equation-key-observation})
we find that the image $\overline{g}'_i \in \overline{S}'$
is not contained in $\mathfrak r'$ for any
$\mathfrak r' \in \text{Ass}_{\overline{S}'}(\overline{N})$.
Hence by Lemma \ref{lemma-invert-universally-injective}
we see that $g'_i : N'_{\mathfrak p'} \to N'_{\mathfrak p'}$ is
$R'_{\mathfrak p'}$-universally injective. In particular
we see that $g'_i : M' \otimes_{R'} N' \to M' \otimes_{R'} N'$ is
injective after localizating at $\mathfrak p'$. The claim
follows because $g_i f_i^{n_i} z' = 0$.

\medskip\noindent
Our claim shows that the annihilator of $z'$ in $R'_{\mathfrak p'}$
contains the elements $f_i^{n_i}$. As $R \to R'$ is \'etale we have
$\mathfrak p'R'_{\mathfrak p'} = \mathfrak pR'_{\mathfrak p'}$
by Algebra, Lemma \ref{algebra-lemma-etale-at-prime}.
Hence the annihilator of $z'$ in $R'_{\mathfrak p'}$ has radical equal to
$\mathfrak p' R_{\mathfrak p'}$ (here we use $z'$ is not zero).
On the other hand
$$
z' \in (M' \otimes_{R'} N')_{\mathfrak p'} =
M'_{\mathfrak p'} \otimes_{R'_{\mathfrak p'}} N'_{\mathfrak p'}
$$
The module $N'_{\mathfrak p'}$ is projective over the local ring
$R'_{\mathfrak p'}$ and hence free
(Algebra, Theorem \ref{algebra-theorem-projective-free-over-local-ring}).
Thus we can find a finite free direct summand $F' \subset N'_{\mathfrak p'}$
such that $z' \in M'_{\mathfrak p'} \otimes_{R'_{\mathfrak p'}} F'$.
If $F'$ has rank $n$, then we deduce that
$\mathfrak p' R'_{\mathfrak p'} \in
\text{WeakAss}_{R'_{\mathfrak p'}}({M'_{\mathfrak p'}}^{\oplus n})$.
This implies
$\mathfrak p'R'_{\mathfrak p'} \in \text{WeakAss}(M'_{\mathfrak p'})$
for example by Algebra, Lemma \ref{algebra-lemma-weakly-ass}.
Then $\mathfrak p' \in \text{WeakAss}_{R'}(M')$
which in turn gives $\mathfrak p \in \text{WeakAss}_R(M)$.
This finishes the proof of the implication
``$\Rightarrow$'' of the equivalence of the lemma.

\medskip\noindent
Assume that $\mathfrak p \in \text{WeakAss}_R(M)$ and
$\overline{\mathfrak q} \in \text{Ass}_{\overline{S}}(\overline{N})$.
We want to show that $\mathfrak q$ is weakly associated to $M \otimes_R N$.
Note that $\overline{\mathfrak q}'$ is a maximal element of
$\text{Ass}_{\overline{S}'}(\overline{N}')$.
This is a consequence of (\ref{equation-key-observation})
and the fact that there are no inclusions among the primes
of $\overline{S}'$ lying over $\overline{\mathfrak q}$
(as fibres of \'etale morphisms are discrete
Morphisms, Lemma \ref{morphisms-lemma-etale-over-field}).
Thus, after replacing $R, S, \mathfrak p, \mathfrak q, M, N$ by
$R', S', \mathfrak p', \mathfrak q', M', N'$
we may assume, in addition to the assumptions of the lemma, that
\begin{enumerate}
\item $\mathfrak p \in \text{WeakAss}_R(M)$,
\item $\overline{\mathfrak q} \in \text{Ass}_{\overline{S}}(\overline{N})$,
\item $N$ is projective as an $R$-module, and
\item $\overline{\mathfrak q}$ is maximal in
$\text{Ass}_{\overline{S}}(\overline{N})$.
\end{enumerate}
There is one more reduction, namely, we may replace
$R, S, M, N$ by their localizations at $\mathfrak p$.
This leads to one more condition, namely,
\begin{enumerate}
\item[(5)] $R$ is a local ring with maximal ideal $\mathfrak p$.
\end{enumerate}
We will finish by showing that (1) -- (5) imply
$\mathfrak q \in \text{WeakAss}(M \otimes_R N)$.

\medskip\noindent
Since $R$ is local and $\mathfrak p \in \text{WeakAss}_R(M)$
we can pick a $y \in M$ whose annihilator $I$ has radical
equal to $\mathfrak p$.
Write $\overline{\mathfrak q} = (\overline{g}_1, \ldots, \overline{g}_n)$
for some $\overline{g}_i \in \overline{S}$. Choose $g_i \in S$
mapping to $\overline{g}_i$.
Then $\mathfrak q = \mathfrak pS + g_1S + \ldots + g_nS$.
Consider the map
$$
\Psi : N/IN \longrightarrow (N/IN)^{\oplus n}, \quad
z \longmapsto (g_1z, \ldots, g_nz).
$$
This is a homomorphism of projective $R/I$-modules.
The local ring $R/I$ is auto-associated
(More on Algebra, Definition \ref{more-algebra-definition-auto-ass})
as $\mathfrak p/I$ is locally nilpotent.
The map $\Psi \otimes \kappa(\mathfrak p)$ is not injective, because
$\overline{\mathfrak q} \in \text{Ass}_{\overline{S}}(\overline{N})$.
Hence More on Algebra, Lemma \ref{more-algebra-lemma-P-fPD-zero}
implies $\Psi$ is not injective. Pick $z \in N/IN$
nonzero in the kernel of $\Psi$. The annihilator $J = \text{Ann}_S(z)$
contains $IS$ and $g_i$ by construction. Thus
$\sqrt{J} \subset S$ contains $\mathfrak q$.
Let $\mathfrak s \subset S$ be a prime minimal over $J$.
Then $\mathfrak q \subset \mathfrak s$,
$\mathfrak s$ lies over $\mathfrak p$, and
$\mathfrak s \in \text{WeakAss}_S(N/IN)$.
The last fact by definition of weakly associated primes.
Apply the ``$\Rightarrow$'' part of the lemma (which we've already proven)
to the ring map $R \to S$ and the modules $R/I$ and $N$
to conclude that
$\overline{\mathfrak s} \in \text{Ass}_{\overline{S}}(\overline{N})$.
Since $\overline{\mathfrak q} \subset \overline{\mathfrak s}$
the maximality of $\overline{\mathfrak q}$, see condition (4) above,
implies that $\overline{\mathfrak q} = \overline{\mathfrak s}$.
This shows that $\mathfrak q = \mathfrak s$ and we conlude
what we want.
\end{proof}